\section{本章训练}

\begin{workedexample}{广度优先搜索与无权最短路}
设从源点 $s$ 出发对一个无权图作广度优先搜索。顶点被放入不同层
\[
    L_k=\{v:d(s,v)=k\}.
\]
当处理 $L_k$ 中的顶点时，每个尚未发现的邻点都会被分配到 $L_{k+1}$。之后不可能再发现更短路径，因为所有距离至多为 $k$ 的顶点都已经处理过。因此，在使用邻接表时，广度优先搜索能以 $O(|V|+|E|)$ 的时间计算最短路距离。
\end{workedexample}

\begin{applicationbox}{图神经网络}
消息传递型图神经网络通过聚合邻居特征来更新顶点表示。经过 $k$ 层后，表示依赖于 $k$ 跳邻域，这与反复乘以由邻接关系导出的算子相平行。图的对称性很重要：聚合应当对邻居顺序不变，而整个模型在顶点重新标号下应当保持一致的变换行为。
\end{applicationbox}

\begin{applicationbox}{网络设计}
生成树刻画无环的最低成本连接；最短路刻画路由；割衡量脆弱性与容量；匹配刻画分配问题；染色刻画无冲突调度。图抽象之所以有用，是因为一旦正确选定顶点和边，同一个定理或算法就能作用于许多系统。
\end{applicationbox}

\begin{chapterexercises}
    \item 证明握手恒等式 $\sum_{v\in V}\deg(v)=2|E|$。
    \item 在一个至少有六个顶点的自选图上运行广度优先搜索，并记录层次和父结点树。
    \item 证明含 $n$ 个顶点的树有 $n-1$ 条边。
    \item 判断奇圈和偶圈的色数。
    \item 证明每个有 $n\ge3$ 个顶点的简单平面图至多有 $3n-6$ 条边。
    \item 在 $G(n,p)$ 中，计算三角形数的期望。
    \item 对 $G(n,c/n)$，用分支过程启发解释为什么 $c=1$ 是临界值。
    \item 给出一个 $G(n,p)$ 不适合建模的网络问题，并指出缺失的结构特征。
\end{chapterexercises}

\begin{selectedhints}
    \item 用两种方式计数关联对 $(v,e)$。
    \item 检查每个非源点是否都在前一层获得一个父结点。
    \item 删除一个叶子并用归纳法，或同时使用连通性和无环性。
    \item 奇圈需要三种颜色；偶圈是二部图，需要两种颜色。
    \item 结合 Euler 公式 $n-m+f=2$ 与 $3f\le2m$。
    \item 可能的三角形有 $\binom n3$ 个，每个以概率 $p^3$ 出现。
    \item 新发现邻居数的期望约为 $c$；均值低于 $1$ 的分支过程会灭绝，高于 $1$ 时可能存活。
    \item 例子包括具有重尾度分布的社交网络、带几何位置的空间网络，或具有高聚类和社群结构的网络。
\end{selectedhints}
